\section{Términos clave y síntesis}

Esta clase reúne en un mismo sistema la representación, el objetivo de generación continua, las modificaciones al Transformer, el curriculum de datos y las aplicaciones multitarea. Si solo se memorizan los nombres de los modelos, es fácil confundir qué componente resuelve cada cuello de botella. Primero se reorganizan los términos en un glosario según el esquema «problema--mecanismo--evidencia» y después se presenta una lista de verificación de extremo a extremo.

\subsection{Glosario de alta densidad}

\begin{center}
\small
\begin{tabularx}{\textwidth}{p{0.18\textwidth} X X}
\toprule
\textbf{Término} & \textbf{Problema que resuelve} & \textbf{Mecanismo central y relación con la clase} \\
\midrule
TAE & Las secuencias de video en píxeles son demasiado largas & Comprime 8 veces el tiempo, la altura y la anchura; unifica imagen y video con un latent continuo. \\
Latent & Qué variable aleatoria se entrega al generador & Conserva la información relevante para la generación y descarta parte de la redundancia de píxeles; la tasa de compresión determina el compromiso entre calidad y costo. \\
Flow Matching & Cómo generar medios a partir de ruido & Aprende un velocity field en tiempo continuo; en la inferencia se integra con un ODE solver. \\
Cross-attention & Cómo leen el texto los tokens de video & El video actúa como query y tres representaciones de texto como key/value; el texto se mantiene como condición limpia. \\
AdaLN & Cómo entra el timestep en cada capa & El embedding del tiempo genera el scale/shift de LayerNorm y modula el mismo bloque. \\
MHA & Interacción espaciotemporal bidireccional & Cada query head tiene sus propias K/V heads; es adecuado para la actualización no autorregresiva de la secuencia completa. \\
GQA & La KV cache autorregresiva es demasiado grande & Varias query heads comparten K/V; Movie Gen T2V no depende de esta ventaja. \\
Context parallelism & 73K tokens no caben en una sola GPU & Particiona según la dimensión de token/context e intercambia la información de atención mediante comunicación colectiva. \\
Curriculum & Entrenar directamente en alta resolución es demasiado caro & A partir de imágenes y video de baja resolución, aumenta progresivamente la resolución, la duración y la complejidad de la tarea. \\
Net Win Rate & Las métricas automáticas son insuficientes & En preferencias humanas por pares, la proporción de victorias de A menos la de B, con su intervalo bootstrap. \\
Post-training & Que el pretraining cubra algo no significa que sea controlable & Usa datos de tarea de alta calidad para mejorar el movimiento, la estética, la edición o la preservación de identidad. \\
World model & Qué tan fuerte es en realidad «parecer razonable» & Hay que distinguir entre plausibility visual, consistencia contrafactual y un modelo causal que permita planificar. \\
\bottomrule
\end{tabularx}
\end{center}

\subsection{Lista de verificación de ingeniería de extremo a extremo}

\begin{enumerate}
  \item \textbf{Representation:} ¿se reportan completos la tasa de compresión espaciotemporal, los latent channels y el patch size?
  \item \textbf{Objective:} ¿se predice noise, velocity o data? ¿Cuáles son el solver y el número de pasos de muestreo?
  \item \textbf{Conditioning:} ¿cómo se reparten las funciones el encoder de texto, la cross-attention, AdaLN y las reference conditions?
  \item \textbf{Context:} ¿a cuántos tokens corresponden la resolución, la tasa de cuadros y la duración? ¿Qué tan grandes son los costos de attention y de activaciones?
  \item \textbf{Data:} ¿son reproducibles los filtros visuales, de movimiento, de contenido, de deduplicación y de captioning?
  \item \textbf{Curriculum:} ¿qué etapas entrenan imágenes, video de baja resolución, video de alta resolución y las ramas de tareas?
  \item \textbf{Evaluation:} ¿se reportan a la vez las muestras fallidas, los prompt fijos, las preferencias humanas, el movimiento y el apego al texto?
  \item \textbf{Deployment:} ¿se hacen públicos los pesos, los ODE steps, la comunicación entre varias GPU, la memoria de video máxima y la evidencia de serving?
\end{enumerate}

\begin{importantbox}{Conclusión central de la clase}
La contribución de Movie Gen puede resumirse en una cadena de sistema: TAE comprime los medios en tokens; Flow Matching define un objetivo de transporte continuo; un Transformer bidireccional de estilo Llama con inicialización aleatoria funciona como backbone escalable; la limpieza de datos y un curriculum progresivo hacen viable el entrenamiento con 73K tokens; y un posentrenamiento independiente extiende el sistema a la edición, la personalización y el audio. La mejora de capacidades es real, pero las interacciones complejas, los videos largos, la evaluación, los reasoning verifier y el serving siguen siendo límites abiertos.
\end{importantbox}
